% ============================================================
% fig_som_resultado.tex — Resultado: mapa semántico organizado
% Uso: % ============================================================
% fig_som_resultado.tex — Resultado: mapa semántico organizado
% Uso: % ============================================================
% fig_som_resultado.tex — Resultado: mapa semántico organizado
% Uso: % ============================================================
% fig_som_resultado.tex — Resultado: mapa semántico organizado
% Uso: \input{figures/fig_som_resultado}
% ============================================================
\begin{figure}[H]
\centering
\begin{tikzpicture}[scale=0.9]

% ─── Antes: datos sin organizar ──────────────────────────────
\node[font=\small\bfseries, gray] at (2.0, 5.2) {Datos de entrada};
\node[font=\tiny, gray]           at (2.0, 4.85) {(alta dimensión, sin orden)};

\foreach \x/\y/\col/\lbl in {
    0.4/4.2/red!60/animal,
    1.1/3.5/red!60/animal,
    0.6/3.1/red!60/animal,
    1.8/4.3/blue!60/vehículo,
    2.5/3.8/blue!60/vehículo,
    2.1/3.0/blue!60/vehículo,
    3.2/4.1/green!60!black/fruta,
    3.7/3.4/green!60!black/fruta,
    3.2/2.8/green!60!black/fruta
}{
    \fill[\col] (\x,\y) circle (3.5pt);
}
\node[font=\tiny, red!70]         at (0.7,  2.65) {animales};
\node[font=\tiny, blue!70]        at (2.3,  2.65) {vehículos};
\node[font=\tiny, green!60!black] at (3.5,  2.65) {frutas};

% ─── Flecha central ──────────────────────────────────────────
\draw[-{Stealth}, very thick, blue!60] (4.4, 3.7) -- (5.4, 3.7);
\node[font=\small, blue!70, align=center] at (4.9, 4.1) {entrenamiento\\SOM};

% ─── Después: mapa 4×4 organizado ───────────────────────────
\node[font=\small\bfseries, gray] at (8.5, 5.2) {Mapa SOM resultante};
\node[font=\tiny, gray]           at (8.5, 4.85) {(2D, topología preservada)};

% Celdas del mapa coloreadas por categoría emergente
\foreach \row/\col/\fillcol in {
    0/0/red!40,   0/1/red!30,   0/2/blue!20,  0/3/blue!40,
    1/0/red!35,   1/1/red!20,   1/2/blue!30,  1/3/blue!35,
    2/0/orange!20,2/1/orange!15,2/2/green!25, 2/3/green!40,
    3/0/orange!30,3/1/orange!25,3/2/green!35, 3/3/green!45
}{
    \fill[\fillcol] (5.8+\col*0.95, 4.4-\row*0.95)
        rectangle (5.8+\col*0.95+0.9, 4.4-\row*0.95-0.9);
    \draw[gray!40] (5.8+\col*0.95, 4.4-\row*0.95)
        rectangle (5.8+\col*0.95+0.9, 4.4-\row*0.95-0.9);
}

% Etiquetas de regiones
\node[font=\footnotesize\bfseries, red!70]
    at (6.75, 3.95) {animales};
\node[font=\footnotesize\bfseries, blue!70]
    at (8.15, 3.95) {vehículos};
\node[font=\footnotesize\bfseries, orange!80]
    at (6.75, 2.1) {mixto};
\node[font=\footnotesize\bfseries, green!60!black]
    at (8.15, 2.1) {frutas};

% Nota de gradiente suave
\draw[dashed, gray!50] (7.25, 4.4) -- (7.25, 0.55);
\draw[dashed, gray!50] (5.8, 2.55) -- (9.6, 2.55);
\node[font=\tiny, gray] at (8.5, 0.35)
    {Regiones similares son vecinas en el mapa (topología preservada)};

\end{tikzpicture}
\caption{Resultado del entrenamiento de un SOM. Los datos de entrada de alta
dimensión (izquierda) se proyectan en una grilla 2D (derecha) donde entradas
similares quedan agrupadas en regiones contiguas. La topología del espacio
original se preserva: las regiones de categorías cercanas semánticamente
(como \textit{animales} y \textit{mixto}) quedan adyacentes en el mapa~\cite{kohonen1990som}.}
\label{fig:som_resultado}
\end{figure}

% ============================================================
\begin{figure}[H]
\centering
\begin{tikzpicture}[scale=0.9]

% ─── Antes: datos sin organizar ──────────────────────────────
\node[font=\small\bfseries, gray] at (2.0, 5.2) {Datos de entrada};
\node[font=\tiny, gray]           at (2.0, 4.85) {(alta dimensión, sin orden)};

\foreach \x/\y/\col/\lbl in {
    0.4/4.2/red!60/animal,
    1.1/3.5/red!60/animal,
    0.6/3.1/red!60/animal,
    1.8/4.3/blue!60/vehículo,
    2.5/3.8/blue!60/vehículo,
    2.1/3.0/blue!60/vehículo,
    3.2/4.1/green!60!black/fruta,
    3.7/3.4/green!60!black/fruta,
    3.2/2.8/green!60!black/fruta
}{
    \fill[\col] (\x,\y) circle (3.5pt);
}
\node[font=\tiny, red!70]         at (0.7,  2.65) {animales};
\node[font=\tiny, blue!70]        at (2.3,  2.65) {vehículos};
\node[font=\tiny, green!60!black] at (3.5,  2.65) {frutas};

% ─── Flecha central ──────────────────────────────────────────
\draw[-{Stealth}, very thick, blue!60] (4.4, 3.7) -- (5.4, 3.7);
\node[font=\small, blue!70, align=center] at (4.9, 4.1) {entrenamiento\\SOM};

% ─── Después: mapa 4×4 organizado ───────────────────────────
\node[font=\small\bfseries, gray] at (8.5, 5.2) {Mapa SOM resultante};
\node[font=\tiny, gray]           at (8.5, 4.85) {(2D, topología preservada)};

% Celdas del mapa coloreadas por categoría emergente
\foreach \row/\col/\fillcol in {
    0/0/red!40,   0/1/red!30,   0/2/blue!20,  0/3/blue!40,
    1/0/red!35,   1/1/red!20,   1/2/blue!30,  1/3/blue!35,
    2/0/orange!20,2/1/orange!15,2/2/green!25, 2/3/green!40,
    3/0/orange!30,3/1/orange!25,3/2/green!35, 3/3/green!45
}{
    \fill[\fillcol] (5.8+\col*0.95, 4.4-\row*0.95)
        rectangle (5.8+\col*0.95+0.9, 4.4-\row*0.95-0.9);
    \draw[gray!40] (5.8+\col*0.95, 4.4-\row*0.95)
        rectangle (5.8+\col*0.95+0.9, 4.4-\row*0.95-0.9);
}

% Etiquetas de regiones
\node[font=\footnotesize\bfseries, red!70]
    at (6.75, 3.95) {animales};
\node[font=\footnotesize\bfseries, blue!70]
    at (8.15, 3.95) {vehículos};
\node[font=\footnotesize\bfseries, orange!80]
    at (6.75, 2.1) {mixto};
\node[font=\footnotesize\bfseries, green!60!black]
    at (8.15, 2.1) {frutas};

% Nota de gradiente suave
\draw[dashed, gray!50] (7.25, 4.4) -- (7.25, 0.55);
\draw[dashed, gray!50] (5.8, 2.55) -- (9.6, 2.55);
\node[font=\tiny, gray] at (8.5, 0.35)
    {Regiones similares son vecinas en el mapa (topología preservada)};

\end{tikzpicture}
\caption{Resultado del entrenamiento de un SOM. Los datos de entrada de alta
dimensión (izquierda) se proyectan en una grilla 2D (derecha) donde entradas
similares quedan agrupadas en regiones contiguas. La topología del espacio
original se preserva: las regiones de categorías cercanas semánticamente
(como \textit{animales} y \textit{mixto}) quedan adyacentes en el mapa~\cite{kohonen1990som}.}
\label{fig:som_resultado}
\end{figure}

% ============================================================
\begin{figure}[H]
\centering
\begin{tikzpicture}[scale=0.9]

% ─── Antes: datos sin organizar ──────────────────────────────
\node[font=\small\bfseries, gray] at (2.0, 5.2) {Datos de entrada};
\node[font=\tiny, gray]           at (2.0, 4.85) {(alta dimensión, sin orden)};

\foreach \x/\y/\col/\lbl in {
    0.4/4.2/red!60/animal,
    1.1/3.5/red!60/animal,
    0.6/3.1/red!60/animal,
    1.8/4.3/blue!60/vehículo,
    2.5/3.8/blue!60/vehículo,
    2.1/3.0/blue!60/vehículo,
    3.2/4.1/green!60!black/fruta,
    3.7/3.4/green!60!black/fruta,
    3.2/2.8/green!60!black/fruta
}{
    \fill[\col] (\x,\y) circle (3.5pt);
}
\node[font=\tiny, red!70]         at (0.7,  2.65) {animales};
\node[font=\tiny, blue!70]        at (2.3,  2.65) {vehículos};
\node[font=\tiny, green!60!black] at (3.5,  2.65) {frutas};

% ─── Flecha central ──────────────────────────────────────────
\draw[-{Stealth}, very thick, blue!60] (4.4, 3.7) -- (5.4, 3.7);
\node[font=\small, blue!70, align=center] at (4.9, 4.1) {entrenamiento\\SOM};

% ─── Después: mapa 4×4 organizado ───────────────────────────
\node[font=\small\bfseries, gray] at (8.5, 5.2) {Mapa SOM resultante};
\node[font=\tiny, gray]           at (8.5, 4.85) {(2D, topología preservada)};

% Celdas del mapa coloreadas por categoría emergente
\foreach \row/\col/\fillcol in {
    0/0/red!40,   0/1/red!30,   0/2/blue!20,  0/3/blue!40,
    1/0/red!35,   1/1/red!20,   1/2/blue!30,  1/3/blue!35,
    2/0/orange!20,2/1/orange!15,2/2/green!25, 2/3/green!40,
    3/0/orange!30,3/1/orange!25,3/2/green!35, 3/3/green!45
}{
    \fill[\fillcol] (5.8+\col*0.95, 4.4-\row*0.95)
        rectangle (5.8+\col*0.95+0.9, 4.4-\row*0.95-0.9);
    \draw[gray!40] (5.8+\col*0.95, 4.4-\row*0.95)
        rectangle (5.8+\col*0.95+0.9, 4.4-\row*0.95-0.9);
}

% Etiquetas de regiones
\node[font=\footnotesize\bfseries, red!70]
    at (6.75, 3.95) {animales};
\node[font=\footnotesize\bfseries, blue!70]
    at (8.15, 3.95) {vehículos};
\node[font=\footnotesize\bfseries, orange!80]
    at (6.75, 2.1) {mixto};
\node[font=\footnotesize\bfseries, green!60!black]
    at (8.15, 2.1) {frutas};

% Nota de gradiente suave
\draw[dashed, gray!50] (7.25, 4.4) -- (7.25, 0.55);
\draw[dashed, gray!50] (5.8, 2.55) -- (9.6, 2.55);
\node[font=\tiny, gray] at (8.5, 0.35)
    {Regiones similares son vecinas en el mapa (topología preservada)};

\end{tikzpicture}
\caption{Resultado del entrenamiento de un SOM. Los datos de entrada de alta
dimensión (izquierda) se proyectan en una grilla 2D (derecha) donde entradas
similares quedan agrupadas en regiones contiguas. La topología del espacio
original se preserva: las regiones de categorías cercanas semánticamente
(como \textit{animales} y \textit{mixto}) quedan adyacentes en el mapa~\cite{kohonen1990som}.}
\label{fig:som_resultado}
\end{figure}

% ============================================================
\begin{figure}[H]
\centering
\begin{tikzpicture}[scale=0.9]

% ─── Antes: datos sin organizar ──────────────────────────────
\node[font=\small\bfseries, gray] at (2.0, 5.2) {Datos de entrada};
\node[font=\tiny, gray]           at (2.0, 4.85) {(alta dimensión, sin orden)};

\foreach \x/\y/\col/\lbl in {
    0.4/4.2/red!60/animal,
    1.1/3.5/red!60/animal,
    0.6/3.1/red!60/animal,
    1.8/4.3/blue!60/vehículo,
    2.5/3.8/blue!60/vehículo,
    2.1/3.0/blue!60/vehículo,
    3.2/4.1/green!60!black/fruta,
    3.7/3.4/green!60!black/fruta,
    3.2/2.8/green!60!black/fruta
}{
    \fill[\col] (\x,\y) circle (3.5pt);
}
\node[font=\tiny, red!70]         at (0.7,  2.65) {animales};
\node[font=\tiny, blue!70]        at (2.3,  2.65) {vehículos};
\node[font=\tiny, green!60!black] at (3.5,  2.65) {frutas};

% ─── Flecha central ──────────────────────────────────────────
\draw[-{Stealth}, very thick, blue!60] (4.4, 3.7) -- (5.4, 3.7);
\node[font=\small, blue!70, align=center] at (4.9, 4.1) {entrenamiento\\SOM};

% ─── Después: mapa 4×4 organizado ───────────────────────────
\node[font=\small\bfseries, gray] at (8.5, 5.2) {Mapa SOM resultante};
\node[font=\tiny, gray]           at (8.5, 4.85) {(2D, topología preservada)};

% Celdas del mapa coloreadas por categoría emergente
\foreach \row/\col/\fillcol in {
    0/0/red!40,   0/1/red!30,   0/2/blue!20,  0/3/blue!40,
    1/0/red!35,   1/1/red!20,   1/2/blue!30,  1/3/blue!35,
    2/0/orange!20,2/1/orange!15,2/2/green!25, 2/3/green!40,
    3/0/orange!30,3/1/orange!25,3/2/green!35, 3/3/green!45
}{
    \fill[\fillcol] (5.8+\col*0.95, 4.4-\row*0.95)
        rectangle (5.8+\col*0.95+0.9, 4.4-\row*0.95-0.9);
    \draw[gray!40] (5.8+\col*0.95, 4.4-\row*0.95)
        rectangle (5.8+\col*0.95+0.9, 4.4-\row*0.95-0.9);
}

% Etiquetas de regiones
\node[font=\footnotesize\bfseries, red!70]
    at (6.75, 3.95) {animales};
\node[font=\footnotesize\bfseries, blue!70]
    at (8.15, 3.95) {vehículos};
\node[font=\footnotesize\bfseries, orange!80]
    at (6.75, 2.1) {mixto};
\node[font=\footnotesize\bfseries, green!60!black]
    at (8.15, 2.1) {frutas};

% Nota de gradiente suave
\draw[dashed, gray!50] (7.25, 4.4) -- (7.25, 0.55);
\draw[dashed, gray!50] (5.8, 2.55) -- (9.6, 2.55);
\node[font=\tiny, gray] at (8.5, 0.35)
    {Regiones similares son vecinas en el mapa (topología preservada)};

\end{tikzpicture}
\caption{Resultado del entrenamiento de un SOM. Los datos de entrada de alta
dimensión (izquierda) se proyectan en una grilla 2D (derecha) donde entradas
similares quedan agrupadas en regiones contiguas. La topología del espacio
original se preserva: las regiones de categorías cercanas semánticamente
(como \textit{animales} y \textit{mixto}) quedan adyacentes en el mapa~\cite{kohonen1990som}.}
\label{fig:som_resultado}
\end{figure}
